% =====================================================================
\section{제안 아키텍처}
\label{sec:arch}

그림~\ref{fig:arch}에는 제안 아키텍처의 운영 경로, 터미널 환경과 오프라인 학습 경로가
나타나 있다. 운영 경로는 새로 통지된 작업을 한 번 검토하고, 독립된 제안·수락
정책과 중앙 확정 절차를 거쳐 배정을 바꾼다. 터미널 환경은 그 결정 이후의 트럭 대기,
추가 이동, 재취급과 선박 유휴비용을 계산한다. 학습 경로는 결정 직전의 상태와 난수를
복제해 관측 행동과 반사실 대안을 비교하며, 이 계산은 운영 경로와 분리된다.

\begin{center}
\begin{minipage}{\textwidth}
\centering
\includegraphics[width=\textwidth]{final_figure/system-architecture.png}
\captionof{figure}{제안 아키텍처의 상세 흐름. (a) 운영 단계는 60초마다 자격 확인, 제안,
수락과 확정을 수행하며, 상호 동의나 용량 조건을 만족하지 못하면 원배정을 유지한다.
(b) 이산사건 환경은 게이트, 21개 블록, 블록당 야드 크레인 2기와 선박 작업을 연결해
운영비 $\Phi$를 계산한다. (c) 오프라인 학습은 같은 상태와 난수 흐름에서 최대 세 세계를
3시간 비교해 제안망과 수락망의 비용 라벨을 각각 만든다.}
\label{fig:arch}
\end{minipage}
\end{center}

\subsection{기호·행동·정책}
\label{sec:arch-notation}
결정 간격은 $\delta=60$초다. 작업 $i$의 현재 배정을 $z_i=(b_i,\tau_i)$로 나타내며,
$b_i$는 블록이고 $\tau_i$는 통지된 도착예정시각이다. 재배치 행동은 다음과 같다.

\begin{equation}
a_i=(b'_i,\Delta_i)\in\mathcal A_i(t),\qquad
z_i(a_i)=(b'_i,\tau_i+\Delta_i).
\end{equation}

원배정 유지는 $a_i^0=(b_i,0)$이다. $b'_i\ne b_i$이면 블록을 바꾸고,
$\Delta_i>0$이면 도착시각을 늦춘다. 두 값을 함께 바꾸는 후보도 허용한다. 다만 반출
작업은 꺼낼 컨테이너의 위치가 이미 정해져 있으므로 블록 재배치 후보에서 제외한다.

작업의 통지시각을 $n_i$, 게이트 진입시각을 $g_i$, 최초 검토 여부를 $r_i$라 한다.
작업은 \emph{통지될 때}가 아니라 \emph{도착이 검토 창 안으로 들어올 때} 자격을 얻는다.
창은 $w_i=\min(W,\ \tau_i-n_i)$ 이고 $W=1800$초다. 리드타임이 허락하는 가장 최신
상태에서 결정하기 위해서다.

\begin{equation}
\mathcal E_t=\{i\mid w_i-\delta<\tau_i-t\le w_i,\ g_i>t,\ r_i=0\}.
\end{equation}

따라서 각 작업은 정확히 한 결정 구간에만 들어오고 다시 검토되지 않는다. 도착보다
$W$ 이상 앞서 통지된 보통 작업은 \textbf{도착 정확히 $W$ 전}에 결정된다. 통지 중앙값보다
약 46분 늦은 시점이며, 정보의 신선함은 여기서 온다. 리드가 그보다 짧은 작업은 통지되는
즉시 검토된다. 도착보다 늦게 통지된 작업은 $w_i=0$ 이라 자격을 한 번도 얻지 못하는데,
\ref{sec:env-demand}절의 리드타임 분포에서 이런 작업이 수요의 14.1\%이며 아키텍처는
그 배정을 건드리지 않는다.

제안 정책은 현재 블록 상태와 작업·행동 특징 $o_i^P$를 이용한다. 후보별 비용 점수를
계산하고 가장 낮은 행동을 선택한다(\ref{sec:arch-models}절).

\begin{equation}
\hat a_i=\arg\min_{a\in\mathcal A_i(t)}
\hat c_{\theta}^{P}(o_i^P,a),\qquad
p_i=\mathbf 1[\hat a_i\ne a_i^0].
\end{equation}

제안이 발생하면 목적 블록 또는 목표 시간대가 수락 주체가 된다. 수락 정책은 제안
메시지와 목적지 상태 $o_i^A$를 이용해 수락과 거절의 예상 비용을 비교한다.

\begin{equation}
q_i=\mathbf 1\!\left[
\hat c_{\psi}^{A}(o_i^A,\hat a_i,\mathrm{accept})
\le
\hat c_{\psi}^{A}(o_i^A,\hat a_i,\mathrm{reject})
\right].
\end{equation}

목표 자원의 남은 용량이 0이면 신경망을 호출하기 전에 거절한다. 두 정책이 모두 동의한
후보의 집합은 $\mathcal J_t=\{i\in\mathcal E_t\mid p_iq_i=1\}$이다.

\subsection{중앙 확정 제약}
\label{sec:arch-commit}
중앙 절차는 새 후보를 만들지 않으며 $\mathcal J_t$ 안에서만 작업을 확정한다. 먼저
동의된 후보를 수락 정책의 예측비용 오름차순으로 정렬한다. 동점이면 작업과 좌표의
사전순을 적용한다. 후보 $i$의 확정 여부를 $x_i\in\{0,1\}$, 목적 블록의 남은 슬롯을
$F_b(t)$, 예약 시간대 $k$의 남은 정원을 $K_k(t)$라 하면 다음 제약을 만족한다.

\begin{equation}
\begin{aligned}
x_i&\le p_iq_i,\qquad
\sum_{i\in\mathcal J_t}\!x_i\mathbf1[b'_i\!=\!b]\le F_b(t),\\
\sum_{i\in\mathcal J_t}\!x_i\mathbf1[k(i)\!=\!k]&\le K_k(t),\qquad
\sum_{t\in\mathcal T}\!\mathbf1[i\in\mathcal E_t]\le1.
\end{aligned}
\end{equation}

블록 용량 제약은 모든 실험에서 적용한다. 시간대 정원 자료가 없는 경우에는
$K_k(t)=\infty$로 둔다. 따라서 이 조건의 결과는 도착시각 조정이 제공할 수 있는 가능성의
범위로 해석한다. 확정된 배정은 결정 시점이 끝날 때 한꺼번에 반영된다. 먼저 확정된
작업이 같은 시점의 다른 후보가 보는 상태를 바꾸지 않도록 하기 위해서다. 이 구조는 두
조건을 보장한다. 제안과 수락이 모두 있는 후보만 확정되며, 각 작업은 한 번만 검토된다.

\subsection{정책모형과 입력}
\label{sec:arch-models}
\ref{sec:arch-notation}절의 $\hat c^{P}_{\theta}$와 $\hat c^{A}_{\psi}$는 비용 신경망이다.
두 망은 행동 확률을 출력하지 않는다. 관측 상태와 후보 하나를 입력받아 스칼라 비용
점수를 출력한다. 이 점수는 절대 비용이 아니라 \textbf{짝 중심 상대 비용}이다. 학습
전에 반사실 짝의 두 행동을 그 평균으로 중심화하므로(\ref{sec:arch-labels}절), 망은 그
후보가 함께 굴린 비교 대상보다 얼마나 비싼지를 예측한다. 평균 차감은 두 행동에
공통인 상수를 없애는 것이므로 후보의 순위, 따라서 식~(1)의 $\arg\min$과 식~(2)의
비교는 달라지지 않는다. 정책은 같은 망으로 모든 후보를 평가한 뒤 학습된 점수의
$\arg\min$을 선택한다. 두 모형은 고정 길이 특징을 처리하는 다층퍼셉트론이며,
모든 블록이 같은 가중치를 공유한다. 제안 정책은 21차원, 수락 정책은 16차원 입력을
사용한다. 두 모형의 은닉층은 모두 64--ReLU--64--ReLU이고, 매개변수는 각각
5{,}633개와 5{,}313개다. 전체 매개변수는 11{,}000개보다 적다. 입력에는 현재 대기,
재고, 남은 크레인 작업량, 공개 도착예정시각, 경로 차이와 목표 시각 주변의 통지 물량이
포함된다. 아직 실현되지 않은 실제 도착·완료시각과 다른 수락 주체의 동시 응답은
제외한다.

그림~\ref{fig:mlp}은 비용 신경망 하나의 구조를 보여 준다. 후보는 고정 길이 벡터
$x\in\mathbb R^{d_{in}}$로 입력된다. 벡터는 세 선형층을 지나며, 앞의 두 층 뒤에는
ReLU 활성함수가 적용된다. 출력은 그 후보의 짝 중심 비용 점수를 나타내는 스칼라 하나다.
이 구조에서는 후보 목록의 길이가 바뀌어도 모형을 다시 학습할 필요가 없다. 또한 모든
블록이 같은 입력 형식과 가중치를 사용한다. 즉 블록별 모형을 따로 배우는 대신, 현재
상태에서 후보를 평가하는 공통 규칙을 학습한다. 두 정책은 입력 차원만 다르며 제안은
21차원, 수락은 16차원이다.

\begin{center}
\begin{minipage}{\textwidth}
\centering
\includegraphics[width=\textwidth]{final_figure/MLP-relu.png}

\vspace{3pt}
$\operatorname{MLP}(x)=W_3\,\operatorname{ReLU}\bigl(W_2\,\operatorname{ReLU}(W_1x+b_1)+b_2\bigr)+b_3$
\vspace{2pt}
\captionof{figure}{비용 신경망 하나의 구조. 입력은 후보 하나에 대한 고정 길이 특징 벡터이고
두 개의 64차원 은닉층에는 ReLU 활성함수를 적용한다. 출력은 그 후보의 짝 중심 비용
점수 하나다. 한 번의 결정은 같은 망으로 후보를 각각 평가한 뒤 $\arg\min$을 취한다.
제안 정책과 수락 정책은 이 구조를 공유하며 입력 차원만 각각 21과 16으로 다르다.}
\label{fig:mlp}
\end{minipage}
\end{center}

\subsection{반사실 비용 라벨}
\label{sec:arch-labels}
각 표본 결정에서 분기 직전의 상태와 난수 흐름을 복제한다. 그런 다음 최대 세 개의
세계를 3시간 동안 실행한다.

\begin{equation}
\begin{aligned}
\phi_i^F&=\Phi_H(\text{관측 행동}),\qquad
\phi_i^P=\Phi_H(\text{제안만 반대}),\\
\phi_i^A&=\Phi_H(\text{수락만 반대}),\qquad H=10{,}800\,\mathrm{s}.
\end{aligned}
\end{equation}

대안 세계는 언제나 정책이 실제로 한 행동을 뒤집는다. 제안 정책이 변경을 제안했으면
대안은 현재 배정을 유지하고, 유지했으면 대안은 현재 가중치에서 가장 싼 후보를
확정한다. 수락 정책의 대안은 수락과 거절을 맞바꾼다. 수락 요청이 없는 결정에서는
제안 비교만 생성한다. 각 정책의 두 행동은 그 짝의 공통 평균을 뺀 뒤 100,000원으로
나누어 학습한다.

\begin{equation}
\bar\phi_i^P=\frac{\phi_i^F+\phi_i^P}{2},\qquad
y_{i,F}^P=\frac{\phi_i^F-\bar\phi_i^P}{100{,}000},\qquad
y_{i,P}^P=\frac{\phi_i^P-\bar\phi_i^P}{100{,}000}.
\end{equation}

즉 목표값은 짝 비용 차이의 $\pm$절반이다. 수락 정책에도 $(\phi_i^F,\phi_i^A)$를 이용한
같은 중심화를 적용한다. 따라서 기준점은 고정된 어떤 행동이 아니라 \textbf{짝 자체}다.
목표값은 그 행동이 같은 세계에서 함께 굴린 단 하나의 대안보다 얼마나 좋았는지·나빴는지를
말한다. 평균을 빼면 두 행동의 순서는 유지하면서 두 세계가 공유하는 비용을 제거할 수
있다.

\subsection{학습 설정과 운영 분리}
\label{sec:arch-train}
하루에 최대 64개 결정을 표본으로 사용한다. 탐색 확률은 0.50에서 0.05까지 선형으로
줄인다. 표본은 결정 단위로 학습 80\%와 검증 20\%로 나눈다. 최적화에는 Adam(학습률
$3\times10^{-4}$)과 Huber 손실($\beta=1.0$)을 사용한다. 일부 반사실 비용 차이가 다른
표본보다 훨씬 크므로, 제곱오차를 사용하면 소수의 표본이 학습을 지배할 수 있다
~\cite{refhuber}. 미니배치는 최대 256행이며 갱신 횟수는
$\max(50,\min(600,4\times\text{학습행 수}))$다. 기울기는 1.0으로 제한한다.

반사실 분기는 여러 CPU 코어에서 병렬로 실행한다. 그날 생성한 라벨은 당일 갱신 뒤
교체한다. 운영 평가에서는 가중치를 고정하고 탐색을 끈다. 반사실 분기도 호출하지 않는다.

따라서 라벨은 고정된 자료가 아니다. 모든 반사실 짝은 그 시점의 정책이 실제로 한
행동에서 갈라져 나오므로, 어떤 상태에 라벨이 붙는지도, 그 상태가 무엇과 비교되는지도
가중치가 바뀌면 함께 움직인다. 한 회차의 라벨은 그 회차의 갱신 뒤 버린다. 이런 점에서
본 학습은 시뮬레이터가 만든 기준선을 쓰는 온폴리시 몬테카를로 이점 추정에 해당한다.
목표값은 지평 끝까지 굴린 완전한 시뮬레이션 결과이며, 부트스트랩도 정책경사도 쓰지
않는다.

두 망을 24회차 학습하는 데 CPU 시간 4.7시간이 걸렸다. 운영 단계에서는 후보마다 작은
신경망을 한 번 통과시킨다. 이 계산량은 함께 수행하는 제약 검사에 비해 작다. 따라서
주요 계산 비용은 반사실 라벨 생성에서 발생하며, 이 작업은 오프라인에서만 수행된다.
